% GENERATED by backtests/spx_candle_imbalance/make_tex.py from backtests/spx_candle_imbalance/results.json
% (the one registered run of preregistration/spx_candle_imbalance.md). Do not edit numbers by hand; re-run the script.
% Macros used: \dollar, \ITM, \tabnotes.

% =====================================================================
\section{Closing-auction imbalances}
\label{sec:imbalance}

From 15:50~ET the New York Stock Exchange publishes the closing-auction imbalance of the stocks it lists, and Nasdaq
disseminates an early imbalance indicator for its own stocks every ten seconds.\footnote{The NYSE moved its
market-on-close cut-off and imbalance publication from 15:45 to 15:50 on 1 April 2019 (SEC Release No.~34-85021).
Nasdaq's early indicator from 15:50 was scheduled for the fourth quarter of 2019 (Release No.~34-87179). Samples of
the exchanges' messages on 30 December 2020 and 6 November 2026, outside the sample, confirmed both schedules.} The
15:50 candle is therefore the first full minute in which the market can react to the published imbalance. That
suggests an explanation of the signal: traders act on the publication during the 15:50 minute, and the index keeps
drifting the same way while the auction absorbs the imbalance. If that channel mattered, the rule should be right
more often on days when its candle points the same way as the published imbalance than on days when it points the
other way. A pre-registered study (plan locked 2026-11-10 19:37 UTC, one registered run) tested this prediction on the rule's
644 signal days of 2021--2026, the days of its holdout.

\paragraph{Design.} For each year the study used the 50 largest common-stock positions of SPY at the previous
year-end, from the fund's N-PORT reports to the SEC; they held between 51\% and 63\% of its net assets. On each signal day it
bought from Databento the imbalance messages that each stock's listing exchange sent between 15:50 and 16:00 (NYSE
Pillar for NYSE-listed stocks, Nasdaq TotalView-ITCH for Nasdaq-listed ones). For each stock it kept the first
closing-auction imbalance record received in the 15:50 minute and summed the signed dollar imbalances,
$D = \sum_i s_i Q_i p_i$, where $s_i$ is $+1$ for a buy imbalance, $-1$ for a sell imbalance and 0 for none, $Q_i$
is the imbalance in shares and $p_i$ the reference price, in exact integer arithmetic. A day with at least 10 usable
stocks and $D \ne 0$ is classified, and it \emph{agrees} if the sign of $D$ is the candle's direction. The outcome
is the rule's own: a hit when the CFD moves from the 15:51 open to the 15:59 close in the candle's direction. The
primary test (H1) compares the hit rates on agreeing and disagreeing days, $d = h_a - h_d$, and passes if $d > 0$
and $t = d/\mathrm{se} \ge 2.0$, with the standard error of two independent proportions. Two secondary tests,
which decide nothing, ask whether the candle agrees with the imbalance on more than half of the days (H2) and
whether the imbalance direction alone predicts the move (H3), with a Holm correction. A block bootstrap gives an
interval for $d$, and the robustness checks of Table~\ref{tab:imbtest} were fixed in the plan, as was the wording
of the verdict. By the plan's power calculation, which treats days as independent, 80\% power at the bar needs a
true $d$ of about 11 to 13 points: the test can detect a large concentration of the edge on agreeing days, not a modest one.

Every signal day was classified. All 50 stocks had a usable first record on each of the 644 days, and the first
closing record of the minute arrived a median of about 0.1 seconds after 15:50:00 on both exchanges (at most
1.1 seconds). The code was reviewed three times by a second reader before the download and the run, and
every defect found was fixed first.

\begin{table}[htbp]
\centering
\caption{The closing-imbalance test: the rule's hit rate on days that agree and disagree with the 15:50 imbalance.}
\label{tab:imbtest}
\tabnotes{Signal days 2021--2026; all 644 were classified. Cells give the value and, in brackets, the number of
days (trades for the option P\&L). Differences of hit rates are in percentage points; $t$ uses the standard error
of two independent proportions, and the move and option P\&L rows use Welch's $t$. The bootstrap draws 20-day blocks
of the 1{,}446 sessions, 10{,}000 of 10{,}000 draws kept. H2 and H3 are exact one-sided binomial tests
against one half (the half benchmark is descriptive), with a Holm correction: H2 rejects,
H3 does not reject. Panel~B was fixed in the plan and cannot change the verdict.
Move: from the 15:51 open to the 15:59 close in the candle's direction, basis points. Option P\&L: \ITM{10} at the
15:51 ask, held to settlement, \$0.09 fee, on the 585 priced trades.}
\begin{tabular}{@{}lrrrr@{}}
\toprule
 & Agreeing days & Disagreeing days & Difference & Statistic \\
\midrule
\multicolumn{5}{@{}l}{\textit{Panel A: registered tests}} \\
H1: hit rate (primary) & 55.9\% (469) & 61.7\% (175) & $-5.9$ & $t = -1.35$ \\
\quad Block-bootstrap 95\% interval & & & \multicolumn{2}{l}{$-14.8$ to $+4.0$} \\
H2: candle agrees with imbalance & \multicolumn{2}{c}{469 of 644 (72.8\%)} & & $p = 3.1\times10^{-32}$ \\
H3: imbalance direction right & \multicolumn{2}{c}{329 of 644 (51.1\%)} & & $p = 0.30$ \\
\midrule
\multicolumn{5}{@{}l}{\textit{Panel B: robustness}} \\
Latest record before 15:51 & 55.6\% (459) & 62.2\% (185) & $-6.6$ & $t = -1.55$ \\
NYSE-listed stocks only & 56.2\% (372) & 59.2\% (272) & $-3.0$ & $t = -0.76$ \\
Nasdaq-listed stocks only & 57.3\% (499) & 57.9\% (145) & $-0.6$ & $t = -0.13$ \\
Adjusted for year and direction & & & $-6.4$ & 12 cells \\
Move in the candle's direction (bp) & $2.41$ (469) & $2.24$ (175) & $+0.16$ & $t = 0.13$ \\
Option P\&L, \$ a trade & $+\dollar80.43$ (423) & $+\dollar33.58$ (162) & $+\dollar46.85$ & $t = 0.79$ \\
\bottomrule
\end{tabular}
\end{table}

\begin{table}[htbp]
\centering
\caption{Signal days that agree and disagree with the 15:50 closing imbalance, by year.}
\label{tab:imbyear}
\tabnotes{Signal days of the rule (a 15:50 body of 60\% or more), 2021--2026. A day agrees when the sign of the
summed dollar imbalance of the 50 largest stocks, from each stock's first closing-imbalance record received in the
15:50 minute, equals the candle's direction. Every signal day was classified, so no day is unclassified. Hit rate:
share of days on which the CFD moved from the 15:51 open to the 15:59 close in the candle's direction, with its
Wilson 95\% interval in percent; bp: the mean of that move in the candle's direction, in basis points. Over all
days the difference in hit rates is $-5.9$ points ($t = -1.35$), with a 20-day block-bootstrap 95\%
interval from $-14.8$ to $+4.0$ points.}
\begin{tabular}{@{}lrrlrrlr@{}}
\toprule
 & & \multicolumn{3}{c}{Agreeing days} & \multicolumn{3}{c}{Disagreeing days} \\
\cmidrule(lr){3-5}\cmidrule(l){6-8}
Year & Days & $n$ & Hit rate (95\% CI) & bp & $n$ & Hit rate (95\% CI) & bp \\
\midrule
2021 & 111 & 92 & 62.0\% (52--71) & $4.67$ & 19 & 78.9\% (57--91) & $5.12$ \\
2022 & 105 & 76 & 57.9\% (47--68) & $4.00$ & 29 & 69.0\% (51--83) & $3.88$ \\
2023 & 113 & 73 & 53.4\% (42--64) & $0.33$ & 40 & 55.0\% (40--69) & $1.18$ \\
2024 & 106 & 69 & 59.4\% (48--70) & $2.59$ & 37 & 43.2\% (29--59) & $-1.90$ \\
2025 & 117 & 91 & 53.8\% (44--64) & $1.39$ & 26 & 69.2\% (50--83) & $3.97$ \\
2026 & 92 & 68 & 47.1\% (36--59) & $0.96$ & 24 & 70.8\% (51--85) & $4.29$ \\
\midrule
All & 644 & 469 & 55.9\% (51--60) & $2.41$ & 175 & 61.7\% (54--69) & $2.24$ \\
\bottomrule
\end{tabular}
\end{table}

\paragraph{Results.} The candle and the published imbalance usually pointed the same way: they agreed on
469 of the 644 days (72.8\%), far more often than half (H2, $p = 3.1\times10^{-32}$).
Agreement did not make the candle more accurate (Table~\ref{tab:imbtest}). It was right on 55.9\% of
agreeing days and on 61.7\% of disagreeing days, so $d = -5.9$ points (se
4.3, $t = -1.35$), and H1 did not pass. The block-bootstrap interval of $d$, from
$-14.8$ to $+4.0$ points, lies below the 11 to 13 points that the test was designed to detect.
The imbalance direction alone was right on 51.1\% of days (H3, $p = 0.30$). The robustness
checks agree: the difference in hit rates stayed negative with the latest record of the minute, with NYSE or Nasdaq
stocks alone and within year and direction. The average move in the candle's direction was about the same on both
kinds of day ($2.41$ against $2.24$ basis points), and the option trade made more on agreeing
days ($+\dollar80.43$ against $+\dollar33.58$ a trade), but not significantly ($t = 0.79$). By year
(Table~\ref{tab:imbyear}), disagreeing days had the higher hit rate in five of the six years, all but 2024.

\paragraph{What the result means.} In the words the plan fixed in advance, the prediction was not confirmed, the
test has limited power, and a failure does not refute the auction explanation. The price may adjust within the 15:50
candle itself and leave little after 15:51, imbalances change after their first publication, the stocks'
imbalances can offset one another in $D$, and 50 stocks are only part of the index. The high agreement is itself an
association, not a cause: the candle and the imbalance can both reflect the same closing orders, such as index
funds' orders known earlier in the day, and because records are timed by when Databento received them, the study
cannot show that the publication came before the price move. What it does show is narrower: on these days the
rule's edge was not concentrated on the days on which its candle agreed with the imbalance first published for the
50 largest stocks. Whether the effect changed when the NYSE moved its publication to 15:50 in 2019, and whether it
depends on the size of the imbalance, were outside the plan.
